\section{Cryptographic Approaches: MPC, FHE, and Secret Sharing}
\label{sec:crypto}

% Mechanism-section template (consistent across §3-§7):
%   Approach · Representative schemes (table) · Performance (reported) ·
%   Limitations / known attacks · graph-RAG applicability.

\subsection{Approach}
% MPC (secret sharing, GC), FHE, and hybrids: compute on data without exposing
% plaintext; the gold standard for confidentiality, bottlenecked by non-linear
% ops (Softmax, GELU) and communication. Map onto §2: TCB = none / non-colluding.
\needswork{describe MPC/FHE/SS for inference and for private retrieval}

\subsection{Representative Schemes}
\begin{table}[t]
  \centering \footnotesize
  \caption{Cryptographic schemes (inference + private retrieval).}
  \label{tab:crypto}
  \begin{tabular}{@{}llll@{}}
    \toprule
    Scheme & Pipeline stage & Primitive & Reported cost \\
    \midrule
    PermLLM     & inference          & SS + permutation & 3\,s/tok (6B, WAN) \\
    Euston      & inference          & non-interactive  & \tq \\
    SecFormer   & inference          & SMPC             & \tq \\
    CipherFormer& inference          & HE (low rounds)  & \tq \\
    CryptoMoE   & MoE inference      & balanced routing & \tq \\
    Caprise     & storage/retrieval  & DCPE             & 2339 vec/s \\
    Fuchsbauer  & storage            & dist.-cmp.-pres. enc & \tq \\
    Compass     & retrieval (search) & encrypted search & \tq \\
    PACMANN     & retrieval (ANN)    & private graph ANN & \tq \\
    V3DB        & retrieval (verify) & ZK proofs        & \tq \\
    \bottomrule
  \end{tabular}
\end{table}
\needswork{populate costs; note which are inference vs retrieval; Compass/PACMANN
also serve graph-RAG retrieval (see \cref{sec:synthesis})}

\subsection{Performance (Reported)}
\needswork{the headline cost reality: MPC inference is 1--2 orders of magnitude
slower (e.g. BumbleBee ~8\,min/token, 7B); retrieval crypto is closer to practical}

\subsection{Limitations and Known Attacks}
\needswork{non-colluding assumption (MPC); non-linear-op overhead; FHE depth}
